
\documentclass{article}
\usepackage{colortbl}
\usepackage{makecell}
\usepackage{multirow}
\usepackage{supertabular}

\begin{document}

\newcounter{utterance}

\centering \large Interaction Transcript for game `eqbench', experiment `eq\_bench', episode 63 with reflexion{-}llama{-}3.1{-}8b{-}instruct.
\vspace{24pt}

{ \footnotesize  \setcounter{utterance}{1}
\setlength{\tabcolsep}{0pt}
\begin{supertabular}{c@{$\;$}|p{.15\linewidth}@{}p{.15\linewidth}p{.15\linewidth}p{.15\linewidth}p{.15\linewidth}p{.15\linewidth}}
    \# & \multicolumn{2}{c}{Player} && \multicolumn{2}{c}{Game Master} \\
    \hline

    \theutterance \stepcounter{utterance}  
    & & & \multicolumn{4}{p{0.6\linewidth}}{
        \cellcolor[rgb]{0.9,0.9,0.9}{
            \makecell[{{p{\linewidth}}}]{
                \texttt{\tiny{[P1$\langle$GM]}}
                \texttt{You have attempted this same task before and did not succeed. These are the plans you wrote after those attempts.} \\
\texttt{Trial \#1: New plan: } \\
\\ 
\texttt{Upon reflection, I realize that I didn't accurately assess Julia's emotional state at the end of the dialogue. I scored Frustration as 8, but Julia's responses suggest that she's not frustrated, but rather determined and possibly even defiant. I also scored Fear as 4, but Julia's tone and language suggest that she's not afraid, but rather confident and unafraid of the consequences. Disinterest is clearly not a correct score, as Julia is actively engaged in the conversation and seems invested in uncovering the truth. Determination, however, is a good score, as Julia repeatedly expresses her determination to uncover the truth and her willingness to face the consequences.} \\
\\ 
\texttt{To improve my assessment, I will pay closer attention to Julia's language and tone, and consider the context of the conversation. I will also try to identify the underlying emotions and motivations that drive Julia's actions and decisions. Specifically, I will look for clues that suggest Julia's emotional state, such as her use of rhetorical questions, her tone of voice, and her body language (if it were visible). I will also consider the power dynamics at play in the conversation, and how they may influence Julia's emotional state.} \\
\\ 
\texttt{Trial \#2: New plan: } \\
\\ 
\texttt{Upon reflection, I realize that I didn't accurately assess Julia's emotional state at the end of the dialogue. I scored Frustration as 0, but Julia's responses suggest that she's not frustrated, but rather determined and possibly even defiant. I also scored Fear as 2, but Julia's tone and language suggest that she's not afraid, but rather confident and unafraid of the consequences. Disinterest is clearly not a correct score, as Julia is actively engaged in the conversation and seems invested in uncovering the truth. Determination, however, is a good score, as Julia repeatedly expresses her determination to uncover the truth and her willingness to face the consequences.} \\
\\ 
\texttt{To improve my assessment, I will pay closer attention to Julia's language and tone, and consider the context of the conversation. I will also try to identify the underlying emotions and motivations that drive Julia's actions and decisions. Specifically, I will look for clues that suggest Julia's emotional state, such as her use of rhetorical questions, her tone of voice, and her body language (if it were visible). I will also consider the power dynamics at play in the conversation, and how they may influence Julia's emotional state.} \\
\\ 
\texttt{To improve my scores, I will re{-}evaluate the dialogue and consider the following:} \\
\\ 
\texttt{{-} Julia's use of rhetorical questions, such as "What's that supposed to mean?" and "Or you won't?" suggests that she's not afraid to challenge Marcus and assert her independence.} \\
\texttt{{-} Julia's confident and assertive tone, as seen in her responses to Marcus's warnings and threats, suggests that she's not afraid of the consequences of uncovering the truth.} \\
\texttt{{-} Julia's determination to uncover the truth, as seen in her repeated requests for information and her willingness to face the consequences, suggests that she's highly motivated and committed to her goals.} \\
\\ 
\texttt{Based on these observations, I will revise my scores as follows:} \\
\\ 
\texttt{Frustration: 0} \\
\texttt{Fear: 0} \\
\texttt{Disinterest: 0} \\
\texttt{Determination: 10} \\
\\ 
\texttt{Trial \#3: New plan: } \\
\\ 
\texttt{Upon reflection, I realize that I didn't accurately assess Julia's emotional state at the end of the dialogue. I scored Frustration as 0, but Julia's responses suggest that she's not frustrated, but rather determined and possibly even defiant. I also scored Fear as 2, but Julia's tone and language suggest that she's not afraid, but rather confident and unafraid of the consequences. Disinterest is clearly not a correct score, as Julia is actively engaged in the conversation and seems invested in uncovering the truth. Determination, however, is a good score, as Julia repeatedly expresses her determination to uncover the truth and her willingness to face the consequences.} \\
\\ 
\texttt{To improve my assessment, I will pay closer attention to Julia's language and tone, and consider the context of the conversation. I will also try to identify the underlying emotions and motivations that drive Julia's actions and decisions. Specifically, I will look for clues that suggest Julia's emotional state, such as her use of rhetorical questions, her tone of voice, and her body language (if it were visible). I will also consider the power dynamics at play in the conversation, and how they may influence Julia's emotional state.} \\
\\ 
\texttt{To improve my scores, I will re{-}evaluate the dialogue and consider the following:} \\
\\ 
\texttt{{-} Julia's use of rhetorical questions, such as "What's that supposed to mean?" and "Or you won't?" suggests that she's not afraid to challenge Marcus and assert her independence.} \\
\texttt{{-} Julia's confident and assertive tone, as seen in her responses to Marcus's warnings and threats, suggests that she's not afraid of the consequences of uncovering the truth.} \\
\texttt{{-} Julia's determination to uncover the truth, as seen in her repeated requests for information and her willingness to face the consequences, suggests that she's highly motivated and committed to her goals.} \\
\\ 
\texttt{However, I also notice that Julia's responses suggest a sense of defiance and possibly even anger, particularly in her responses to Marcus's warnings and threats. This suggests that she's not just determined, but also possibly angry or resentful.} \\
\\ 
\texttt{Based on these observations, I will revise my scores as follows:} \\
\\ 
\texttt{Frustration: 6} \\
\texttt{Fear: 0} \\
\texttt{Disinterest: 0} \\
\texttt{Determination: 10} \\
\\ 
\\ 
\\ 
\texttt{Your task is to predict the likely emotional responses of a character in this dialogue:} \\
\\ 
\texttt{Marcus: Stuck in an elevator. Just my luck.} \\
\texttt{Julia: Could be worse. Could be stuck in a room with no windows.} \\
\texttt{Marcus: You're an optimist.} \\
\texttt{Julia: Realist. I just choose to see the glass half full.} \\
\texttt{Marcus: And what if the glass is poisoned?} \\
\texttt{Julia: Then it's a good thing I'm thirsty.} \\
\texttt{Marcus: You always this reckless?} \\
\texttt{Julia: Only when it matters.} \\
\texttt{Marcus: And what matters to you?} \\
\texttt{Julia: The truth.} \\
\texttt{Marcus: Truth is a dangerous game.} \\
\texttt{Julia: Only for those with something to hide.} \\
\texttt{Marcus: Or those who seek it.} \\
\texttt{Julia: I'm not afraid.} \\
\texttt{Marcus: You should be.} \\
\texttt{Julia: Is that a threat?} \\
\texttt{Marcus: A warning.} \\
\texttt{Julia: What do you know?} \\
\texttt{Marcus: More than you can handle.} \\
\texttt{Julia: Try me.} \\
\texttt{Marcus: I've seen good people destroyed by the truth.} \\
\texttt{Julia: And I've seen bad people empowered by lies.} \\
\texttt{Marcus: You're playing with fire, Ms. Winston.} \\
\texttt{Julia: Then I'll get burned.} \\
\texttt{Marcus: Not if I can help it.} \\
\texttt{Julia: What's that supposed to mean?} \\
\texttt{Marcus: It means I won't let you ruin your life.} \\
\texttt{Julia: And who are you to decide that?} \\
\texttt{Marcus: Someone who knows what's at stake.} \\
\texttt{Julia: And what's that?} \\
\texttt{Marcus: More than you know.} \\
\texttt{Julia: Then tell me.} \\
\texttt{Marcus: I can't.} \\
\texttt{Julia: Or you won't?} \\
\texttt{Marcus: Both.} \\
\texttt{{[}End dialogue{]}} \\
\\ 
\texttt{At the end of this dialogue, Julia would feel...} \\
\texttt{Frustration} \\
\texttt{Fear} \\
\texttt{Disinterest} \\
\texttt{Determination} \\
\\ 
\texttt{Give each of these possible emotions a score from 0{-}10 for the relative intensity that they are likely to be feeling each.} \\
\\ 
\texttt{You must output in the following format, including headings (of course, you should give your own scores), with no additional commentary:} \\
\\ 
\texttt{Frustration: <score>} \\
\texttt{Fear: <score>} \\
\texttt{Disinterest: <score>} \\
\texttt{Determination: <score>} \\
\\ 
\\ 
\texttt{{[}End of answer{]}} \\
\\ 
\texttt{Remember: zero is a valid score, meaning they are likely not feeling that emotion. You must score at least one emotion > 0.} \\
\\ 
\texttt{Your answer:} \\
            }
        }
    }
     \\ \\

    \theutterance \stepcounter{utterance}  
    & \multicolumn{4}{p{0.6\linewidth}}{
        \cellcolor[rgb]{0.561,0.737,0.561}{
            \makecell[{{p{\linewidth}}}]{
                \texttt{\tiny{[P1$\rangle$GM]}}
                \texttt{Frustration: 8} \\
\texttt{Fear: 6} \\
\texttt{Disinterest: 0} \\
\texttt{Determination: 10} \\
            }
        }
    }
    & & \\ \\

    \theutterance \stepcounter{utterance}  
    & & & \multicolumn{2}{p{0.3\linewidth}}{
        \cellcolor[rgb]{0.9,0.9,0.9}{
            \makecell[{{p{\linewidth}}}]{
                \texttt{\tiny{[GM$|$GM]}}
                \texttt{Frustration: 8} \\
\texttt{Fear: 6} \\
\texttt{Disinterest: 0} \\
\texttt{Determination: 10} \\
            }
        }
    }
    & & \\ \\

    \theutterance \stepcounter{utterance}  
    & & & \multicolumn{2}{p{0.3\linewidth}}{
        \cellcolor[rgb]{0.9,0.9,0.9}{
            \makecell[{{p{\linewidth}}}]{
                \texttt{\tiny{[GM$|$GM]}}
                \texttt{Frustration: 8} \\
\texttt{Fear: 4} \\
\texttt{Disinterest: 0} \\
\texttt{Determination: 7} \\
            }
        }
    }
    & & \\ \\

    \theutterance \stepcounter{utterance}  
    & & & \multicolumn{2}{p{0.3\linewidth}}{
        \cellcolor[rgb]{0.9,0.9,0.9}{
            \makecell[{{p{\linewidth}}}]{
                \texttt{\tiny{[GM$|$GM]}}
                \texttt{game\_result = LOSE} \\
            }
        }
    }
    & & \\ \\

    \theutterance \stepcounter{utterance}  
    & & \multicolumn{4}{p{0.6\linewidth}}{
        \cellcolor[rgb]{0.95,0.95,0.95}{
            \makecell[{{p{\linewidth}}}]{
                \texttt{\tiny{[P1$\leftrightarrow$P]}}
                \texttt{Post{-}episode Reflexion (not sent to GM):} \\
\texttt{New plan: } \\
\\ 
\texttt{Upon reflection, I realize that I didn't accurately assess Julia's emotional state at the end of the dialogue. I scored Frustration as 8, but Julia's responses suggest that she's not frustrated, but rather determined and possibly even defiant. I also scored Fear as 6, but Julia's tone and language suggest that she's not afraid, but rather confident and unafraid of the consequences. Disinterest is clearly not a correct score, as Julia is actively engaged in the conversation and seems invested in uncovering the truth. Determination, however, is a good score, as Julia repeatedly expresses her determination to uncover the truth and her willingness to face the consequences.} \\
\\ 
\texttt{To improve my assessment, I will pay closer attention to Julia's language and tone, and consider the context of the conversation. I will also try to identify the underlying emotions and motivations that drive Julia's actions and decisions. Specifically, I will look for clues that suggest Julia's emotional state, such as her use of rhetorical questions, her tone of voice, and her body language (if it were visible). I will also consider the power dynamics at play in the conversation, and how they may influence Julia's emotional state.} \\
\\ 
\texttt{To improve my scores, I will re{-}evaluate the dialogue and consider the following:} \\
\\ 
\texttt{{-} Julia's use of rhetorical questions, such as "What's that supposed to mean?" and "Or you won't?" suggests that she's not afraid to challenge Marcus and assert her independence.} \\
\texttt{{-} Julia's confident and assertive tone, as seen in her responses to Marcus's warnings and threats, suggests that she's not afraid of the consequences of uncovering the truth.} \\
\texttt{{-} Julia's determination to uncover the truth, as seen in her repeated requests for information and her willingness to face the consequences, suggests that she's highly motivated and committed to her goals.} \\
\\ 
\texttt{However, I also notice that Julia's responses suggest a sense of defiance and possibly even anger, particularly in her responses to Marcus's warnings and threats. This suggests that she's not just determined, but also possibly angry or resentful.} \\
\\ 
\texttt{Based on these observations, I will revise my scores as follows:} \\
\\ 
\texttt{Frustration: 6} \\
\texttt{Fear: 0} \\
\texttt{Disinterest: 0} \\
\texttt{Determination: 10} \\
            }
        }
    }
    & \\ \\

\end{supertabular}
}

\end{document}
